\clearpage\pagestyle{otMain}\fontsize{11}{15}\selectfont
\corrigesChapitre{5}{Fonctions et représentations}
\begin{blocLivre}{cor-5-1}
\begin{corrige}{1}{Calculer des images}
\demarche{Substituer une valeur à toutes les occurrences de la variable.}
On remplace chaque occurrence de $x$ en respectant les parenthèses. $f(0)=\res1$ ; $f(2)=2\times4-6+1=\res3$ ; $f(-1)=2\times1+3+1=\res6$. En particulier $-3(-1)=+3$, ce qui explique le signe dans le dernier calcul.
\end{corrige}
\end{blocLivre}
\begin{blocLivre}{cor-5-2}
\begin{corrige}{2}{Lire un tableau}
\demarche{Lire les colonnes dans le sens correspondant à la question.}
La colonne d’abscisse $-1$ donne $\res{f(-1)=1}$. Les colonnes où la deuxième ligne vaut 4 ont pour abscisses $-2$ et 2. Ce sont les deux antécédents visibles dans le tableau. Sans autre information sur la fonction, on ne peut pas exclure d’autres antécédents hors de ces cinq valeurs.
\end{corrige}
\end{blocLivre}
\begin{blocLivre}{cor-5-3}
\begin{corrige}{3}{Antécédent d’une fonction affine}
\demarche{Traduire le mot antécédent par une équation.}
On cherche $x$ tel que $3x-5=7$. Cela donne $3x=12$, donc $\res{x=4}$. Comme le coefficient 3 est non nul, il n’y a qu’un antécédent. Contrôle : $f(4)=12-5=7$.
\end{corrige}
\end{blocLivre}
\begin{blocLivre}{cor-5-4}
\begin{corrige}{4}{Appartenir à une courbe}
\demarche{Tester l’égalité y = f(x) point par point.}
On compare l’ordonnée annoncée à l’image de l’abscisse. $f(1)=3$, donc $A$ appartient à la courbe ; $f(-2)=6$, donc $B$ aussi. Mais $f(3)=11\ne10$, donc $C$ n’y appartient pas. Une proximité sur un dessin ne remplacerait pas ces égalités exactes.
\label{fig-cor-5-4}
\figureGeo[corrige]{fonction-appartenance}
\end{corrige}
\end{blocLivre}
\begin{blocLivre}{cor-5-5}
\begin{corrige}{5}{Le domaine d’un quotient}
\demarche{Rechercher les zéros du dénominateur avant de calculer.}
Il faut $2x-6\ne0$, donc $x\ne3$. Le domaine est $\res{D_g=\mathbb R\setminus\{3\}}$. À $x=4$, le dénominateur vaut 2 et $\res{g(4)=5/2}$. Il n’existe aucune image de 3 par cette formule.
\end{corrige}
\end{blocLivre}
\begin{blocLivre}{cor-5-6}
\begin{corrige}{6}{Un domaine avec deux contraintes}
\demarche{Prendre l’intersection des conditions de chaque opération.}
La racine impose $x+1\ge0$, soit $x\ge-1$. Le dénominateur impose $x\ne2$. On réunit ces deux contraintes : $\res{D_h=[-1;2[\cup]2;+\infty[}$. La valeur $-1$ est autorisée puisque la racine de 0 existe et que le dénominateur vaut alors $-3$.
\end{corrige}
\end{blocLivre}
\begin{blocLivre}{cor-5-7}
\begin{corrige}{7}{Plusieurs antécédents}
\demarche{Ramener chaque recherche à une équation de carré.}
$(x-1)^2=9$ donne $x-1=-3$ ou 3, donc $\res{x=-2\text{ ou }4}$. Pour 0, il faut $x-1=0$, d’où l’unique antécédent 1. Pour $-1$, il n’y a aucun antécédent réel puisqu’un carré ne peut être négatif.
\label{fig-cor-5-7}
\figureGeo[corrige]{antecedents-parabole}
\end{corrige}
\end{blocLivre}
\begin{blocLivre}{cor-5-8}
\begin{corrige}{8}{Formule et programme de calcul}
\demarche{Transformer les étapes du programme en une seule expression.}
En notant $x$ l’entrée, le résultat est $(4x-6)/2=\res{2x-3}$. Pour produire 11, on résout $2x-3=11$, donc $\res{x=7}$. Faire fonctionner le programme avec 7 donne $(28-6)/2=11$. L’ordre des opérations est décrit par les parenthèses.
\end{corrige}
\end{blocLivre}
\begin{blocLivre}{cor-5-9}
\begin{corrige}{9}{Intersections de deux courbes}
\demarche{Résoudre l’égalité, puis étudier le signe de la différence.}
Les intersections satisfont $x^2-x-2=0$, soit $(x-2)(x+1)=0$. Elles sont $\res{(-1;1)\text{ et }(2;4)}$. Le produit est négatif entre les racines ; ainsi $\res{f(x)\le g(x)\iff x\in[-1;2]}$. Le graphique confirme que la parabole est alors sous la droite.
\label{fig-cor-5-9}
\figureGeo[corrige]{intersection-fonctions}
\end{corrige}
\end{blocLivre}
\begin{blocLivre}{cor-5-10}
\begin{corrige}{10}{Une lecture graphique précise}
\demarche{Explorer tous les morceaux de la courbe, y compris les extrémités.}
Le sommet d’abscisse 0 donne $\res{f(0)=3}$. La droite $y=1$ coupe la ligne en $x=-2$, $x=1$ et $x=4$. Sur le segment central, l’ordonnée descend de 3 à $-1$ en deux unités : elle vaut bien 1 à $x=1$. Les antécédents sont donc $\res{-2,1,4}$.
\label{fig-cor-5-10}
\figureGeo[corrige]{fonction-morceaux}
\end{corrige}
\end{blocLivre}
\begin{blocLivre}{cor-5-11}
\begin{corrige}{11}{Un tarif avec plafond}
\demarche{Décrire séparément la partie proportionnelle et le plafond.}
Jusqu’à 5 h, $C(t)=3t$ ; ensuite $C(t)=15$. Le coût de 12 euros est obtenu seulement pour $t=4$. Le coût de 15 euros est atteint pour toutes les durées $\res{t\in[5;8]}$. Le modèle est continu si les fractions d’heure sont facturées proportionnellement, ce que l’énoncé suppose ici.
\label{fig-cor-5-11}
\figureGeo[corrige]{tarif-plafond}
\end{corrige}
\end{blocLivre}
\begin{blocLivre}{cor-5-12}
\begin{corrige}{12}{Une quantité discrète}
\demarche{Préciser le domaine avant d’accepter une solution algébrique.}
Chaque rangée ajoute 6 : $\res{u(n)=18+6n}$. Obtenir 75 demanderait $n=(75-18)/6=9{,}5$, qui n’est pas entier. Ce total est donc impossible dans ce modèle. Les totaux voisins sont $u(9)=72$ et $u(10)=78$. Le domaine discret est essentiel pour conclure.
\end{corrige}
\end{blocLivre}
\begin{blocLivre}{cor-5-13}
\begin{corrige}{13}{Deux formules, un même domaine ?}
\demarche{Distinguer égalité des valeurs et égalité des fonctions.}
$f$ est définie pour $x\ne1$ ; $g$ l’est pour tout réel. Sur le domaine de $f$, on simplifie $(x-1)(x+1)/(x-1)$, donc $f(x)=g(x)$. La courbe de $f$ est la droite de $g$ privée du point $(1;2)$. Les formules coïncident là où toutes deux sont définies, mais les fonctions n’ont pas le même domaine.
\end{corrige}
\end{blocLivre}
\begin{blocLivre}{cor-5-14}
\begin{corrige}{14}{Concevoir un modèle de volume}
\demarche{Associer une formule géométrique au domaine physique de la variable.}
Le volume d’un cylindre est aire de base fois hauteur : $\res{V(h)=4\pi h}$, pour $\res{h\in[0;3]}$. L’équation $4\pi h=6\pi$ donne $\res{h=1{,}5\text{ m}}$. Cette hauteur appartient au domaine. La capacité totale est $12\pi$ m$^3$, donc le volume demandé est la moitié de la capacité.
\label{fig-cor-5-14}
\figureGeo[corrige]{cylindre}
\end{corrige}
\end{blocLivre}
\begin{blocLivre}{cor-5-15}
\begin{corrige}{15}{Une conjecture issue d’un tableau}
\demarche{Transformer une observation numérique en preuve par un carré.}
On réécrit $f(x)=(x-2)^2-1$. Un carré est positif ou nul, donc $f(x)\ge-1$ pour tout réel. L’égalité exige $(x-2)^2=0$, soit $\res{x=2}$. Le tableau proposait une piste ; la transformation algébrique prouve que l’inégalité reste vraie entre les valeurs testées et au-delà.
\end{corrige}
\end{blocLivre}
